% Required packages in the main paper:
% \usepackage{amsmath,amssymb,bm}

\section{Method}
\label{sec:method}

\subsection{Overview}
\label{sec:method_overview}

We consider multi-view 3D reconstruction from a sequence of LDR images and
camera-aligned event streams under severe overexposure. Let
$\bm{I}_{i}^{e}$ denote the RGB observation of view $i$ at exposure $e$,
$\bm{V}_{i}$ its event voxel, and $\bm{K}_{i}$ the camera intrinsics. Our goal
is to estimate depth $D_i$, point map $\bm{P}_i$, and camera pose $\bm{T}_i$.
At inference time, only one LDR exposure and its corresponding event stream
are required.

Overexposure removes texture gradients, boundaries, and local shading cues
from RGB images. Events provide complementary temporal contrast, but their
responses are ambiguous: highlights, reflections, unrelated illumination
changes, and sensor noise may trigger events without corresponding to scene
geometry. Event measurements are also spatially sparse, making aggressive
input filtering likely to remove weak but useful boundary evidence.

Our method addresses these issues using three components. First,
\emph{paired-exposure token alignment} learns image features that are stable
across different LDR levels of the same geometry. Second,
\emph{geometry-detail supervision} explicitly allocates optimization gradients
to boundaries and local surface-normal variations. Third,
\emph{exposure-consistent event reliability} distills event support, geometric
detail, and token agreement into a dense reliability map. The complete event
voxel is preserved during encoding; reliability controls only the geometric
residual written to the RGB prediction.

\subsection{Event Representation and Coarse Geometry}
\label{sec:event_and_coarse}

For each view interval, events are divided into $B$ temporal bins and positive
and negative polarities are accumulated separately:
\begin{equation}
    \bm{V}_{i}\in\mathbb{R}^{2B\times H\times W},\qquad
    \bm{V}_{i}=
    [V_{i,1}^{+},\ldots,V_{i,B}^{+},V_{i,1}^{-},\ldots,V_{i,B}^{-}].
    \label{eq:event_voxel}
\end{equation}
Event counts are compressed before encoding:
\begin{equation}
    \widehat{V}_{i,c}(\bm{u})=
    \frac{\log\!\left(1+\operatorname{clip}
    (V_{i,c}(\bm{u}),0,c_{\max})\right)}
    {\log(1+c_{\max})},
    \label{eq:event_normalization}
\end{equation}
where $\bm{u}$ denotes a pixel and $c$ indexes event channels.

An RGB geometry backbone first predicts globally stable coarse geometry:
\begin{equation}
    \{D_i^{c},\bm{P}_i^{c},\bm{T}_i\}_{i=1}^{N}
    =\mathcal{G}_{\mathrm{rgb}}
    (\{\bm{I}_{i}^{e_i}\}_{i=1}^{N}).
    \label{eq:coarse_geometry}
\end{equation}
The event branch is subsequently used for spatially localized geometry
refinement rather than global pose or scale estimation.

\subsection{Paired-Exposure Token Alignment}
\label{sec:paired_token_alignment}

During training, we sample two LDR renderings $\bm{I}_{i}^{a}$ and
$\bm{I}_{i}^{b}$ of the same view, pose, geometry, and event interval. The
appearance changes with exposure, while the underlying scene geometry and
event observation remain fixed. Let $\bm{Z}_{i,j}^{e}$ be patch token $j$ at
exposure $e\in\{a,b\}$. We insert a bounded residual adapter into selected
image-aggregator layers:
\begin{equation}
    \widetilde{\bm{Z}}_{i,j}^{e}
    =\bm{Z}_{i,j}^{e}
    +\alpha_{z}\tanh\!\left[
    \bm{W}_{2}\,\phi\!\left(
    \bm{W}_{1}\operatorname{LN}(\bm{Z}_{i,j}^{e})
    \right)\right],
    \label{eq:token_adapter}
\end{equation}
where $\phi$ is a nonlinear activation and $\alpha_z$ bounds the adaptation.
The final projection $\bm{W}_{2}$ is zero-initialized so that optimization
starts from the pretrained geometry representation.

We select the less saturated exposure as the alignment anchor using
\begin{equation}
    q(\bm{I})=
    1-\operatorname{Sat}(\bm{I})
    +\eta\operatorname{Contrast}(\bm{I}).
    \label{eq:anchor_quality}
\end{equation}
Assuming exposure $a$ is selected as the anchor, token alignment is supervised
by
\begin{equation}
    \mathcal{L}_{\mathrm{tok}}=
    \frac{1}{NJ}\sum_{i=1}^{N}\sum_{j=1}^{J}
    \operatorname{SmoothL1}\!\left(
    \operatorname{LN}(\widetilde{\bm{Z}}_{i,j}^{b}),
    \operatorname{sg}\!\left[
    \operatorname{LN}(\widetilde{\bm{Z}}_{i,j}^{a})
    \right]\right),
    \label{eq:token_alignment_loss}
\end{equation}
where $\operatorname{sg}$ denotes stop-gradient. Feature-level alignment avoids
forcing two dense depth maps toward an over-smoothed average while removing
exposure-specific variation before geometry decoding.

The aligned tokens also provide a spatial exposure-agreement cue. We compute
\begin{equation}
    c_{i,j}=
    \frac{\langle
    \widetilde{\bm{Z}}_{i,j}^{a},
    \widetilde{\bm{Z}}_{i,j}^{b}\rangle}
    {\|\widetilde{\bm{Z}}_{i,j}^{a}\|_2
     \|\widetilde{\bm{Z}}_{i,j}^{b}\|_2},
\end{equation}
followed by threshold normalization and spatial upsampling:
\begin{equation}
    C_i(\bm{u})=
    \operatorname{Up}\!\left(
    \left\{
    \operatorname{clip}\!\left(
    \frac{c_{i,j}-\tau_c}{1-\tau_c},0,1
    \right)
    \right\}_{j=1}^{J}
    \right)(\bm{u}).
    \label{eq:exposure_agreement}
\end{equation}
Here $C_i$ measures exposure stability and is not treated as geometric
reliability by itself.

\subsection{Geometry-Detail Supervision}
\label{sec:geometry_detail}

Standard depth and point losses are dominated by large, low-frequency
surfaces and may obtain favorable average errors while missing thin boundaries,
small concavities, and local normal transitions. We therefore construct a
geometry-detail cue from ground-truth depth $D_i^{\ast}$. Surface normals are
computed using camera intrinsics:
\begin{equation}
    \bm{N}_{i}^{\ast}=\Pi(D_i^{\ast},\bm{K}_i),
    \label{eq:gt_normal}
\end{equation}
where $\Pi$ denotes depth-to-normal conversion. The raw detail magnitude is
\begin{equation}
    G_{i}^{\mathrm{raw}}(\bm{u})=
    \frac{1}{2}\left(
    \|\nabla_x\bm{N}_{i}^{\ast}(\bm{u})\|_2+
    \|\nabla_y\bm{N}_{i}^{\ast}(\bm{u})\|_2
    \right).
\end{equation}
We apply robust quantile normalization and a square-root transform:
\begin{equation}
    G_i(\bm{u})=
    \sqrt{
    \operatorname{clip}\!\left(
    \frac{G_{i}^{\mathrm{raw}}(\bm{u})}
    {Q_{0.95}(G_{i}^{\mathrm{raw}})},0,1
    \right)}.
    \label{eq:geometry_detail_map}
\end{equation}
The square root increases the relative contribution of weak but meaningful
surface-normal transitions.

Let $M_i$ be the valid-depth mask and let the event-support map be
\begin{equation}
    S_i(\bm{u})=
    \operatorname{Dilate}\!\left(
    \mathbf{1}\!\left[
    \sum_{c=1}^{2B}|V_{i,c}(\bm{u})|>0
    \right]\right).
    \label{eq:event_support_map}
\end{equation}
Events are used as an emphasis cue rather than as a geometry label:
\begin{equation}
    W_i^{\mathrm{det}}(\bm{u})=
    M_i(\bm{u})G_i(\bm{u})
    \left[1+\beta S_i(\bm{u})\right].
    \label{eq:detail_weight}
\end{equation}
Thus, an event alone does not imply a geometric discontinuity, but a GT
geometry detail supported by events receives stronger supervision.

Let $\widehat{\bm{N}}_i=\Pi(D_i,\bm{K}_i)$ denote the predicted normal map.
For a per-pixel quantity $X_i$, define the weighted mean
\begin{equation}
    \mathbb{E}_{W}[X]=
    \frac{\sum_{i,\bm{u}}W_i(\bm{u})X_i(\bm{u})}
    {\sum_{i,\bm{u}}W_i(\bm{u})+\epsilon}.
\end{equation}
The geometry-detail loss combines normal orientation, normal gradients, and
high-frequency normal residuals:
\begin{align}
    \mathcal{L}_{\mathrm{ori}}
    &=\mathbb{E}_{W^{\mathrm{det}}}
    \left[1-
    \left\langle
    \widehat{\bm{N}}_i,\bm{N}_i^{\ast}
    \right\rangle\right],\\
    \mathcal{L}_{\mathrm{grad}}
    &=\mathbb{E}_{W^{\mathrm{det}}}
    \left[\left\|
    \nabla\widehat{\bm{N}}_i-
    \nabla\bm{N}_i^{\ast}
    \right\|_1\right],\\
    \mathcal{L}_{\mathrm{hf}}
    &=\mathbb{E}_{W^{\mathrm{det}}}
    \left[\left\|
    \operatorname{HF}(\widehat{\bm{N}}_i)-
    \operatorname{HF}(\bm{N}_i^{\ast})
    \right\|_1\right],
\end{align}
and
\begin{equation}
    \boxed{
    \mathcal{L}_{\mathrm{det}}=
    \lambda_{\mathrm{ori}}\mathcal{L}_{\mathrm{ori}}+
    \lambda_{\mathrm{grad}}\mathcal{L}_{\mathrm{grad}}+
    \lambda_{\mathrm{hf}}\mathcal{L}_{\mathrm{hf}}
    }.
    \label{eq:detail_loss}
\end{equation}

\subsection{Exposure-Consistent Event Reliability}
\label{sec:event_reliability}

Event support $S_i$ indicates where events exist but cannot reject highlights
or reflections. Geometry detail $G_i$ identifies real surface variation but
does not indicate whether it is observed by events. Token agreement $C_i$
measures exposure stability but may also respond to stable texture. We combine
the three complementary cues into a selective soft target:
\begin{equation}
    R_i^{\ast}(\bm{u})=
    M_i(\bm{u})S_i(\bm{u})G_i(\bm{u})C_i(\bm{u}).
    \label{eq:reliability_target}
\end{equation}
A small dilation accounts for spatial misalignment between event boundaries,
normal gradients, and patch-token locations:
\begin{equation}
    \overline{R}_i^{\ast}=
    \operatorname{Dilate}(R_i^{\ast}).
\end{equation}
The multiplicative construction assigns high reliability only when event
evidence, geometric variation, and exposure stability are simultaneously
present.

A lightweight ReliabilityNet predicts reliability from one RGB observation
and the complete event voxel:
\begin{equation}
    \widehat{R}_i^{e}=
    \sigma\!\left(
    \mathcal{R}_{\theta}
    ([\widehat{\bm{V}}_i,\bm{I}_i^{e}])
    \right),\qquad e\in\{a,b\}.
    \label{eq:reliability_prediction}
\end{equation}
The regression objective is
\begin{equation}
    \mathcal{L}_{\mathrm{rel}}=
    \frac{1}{2}\sum_{e\in\{a,b\}}
    \operatorname{WBCE}
    (\widehat{R}_i^{e},\overline{R}_i^{\ast}).
    \label{eq:reliability_bce}
\end{equation}
Because paired observations share geometry and events, their reliability maps
should be exposure invariant:
\begin{equation}
    \mathcal{L}_{\mathrm{pair}}=
    \frac{\sum_{i,\bm{u}}M_i(\bm{u})
    |\widehat{R}_i^{a}(\bm{u})-
     \widehat{R}_i^{b}(\bm{u})|}
    {\sum_{i,\bm{u}}M_i(\bm{u})+\epsilon}.
    \label{eq:reliability_pair}
\end{equation}

We further separate positive and negative locations using a ranking loss. Let
\begin{equation}
    \Omega_i^{+}=\{\bm{u}\mid
    \overline{R}_i^{\ast}(\bm{u})\geq\tau_{+}\},\qquad
    \Omega_i^{-}=\{\bm{u}\mid
    \overline{R}_i^{\ast}(\bm{u})\leq\tau_{-}\}.
\end{equation}
Then
\begin{equation}
    \mathcal{L}_{\mathrm{rank}}=
    \frac{1}{2N}\sum_{i=1}^{N}\sum_{e\in\{a,b\}}
    \max\!\left(
    0,m_r-\mu_{\Omega_i^{+}}(\widehat{R}_i^{e})
    +\mu_{\Omega_i^{-}}(\widehat{R}_i^{e})
    \right),
    \label{eq:reliability_rank}
\end{equation}
where $m_r$ is the ranking margin. The complete reliability loss is
\begin{equation}
    \boxed{
    \mathcal{L}_{R}=
    \mathcal{L}_{\mathrm{rel}}+
    \lambda_{\mathrm{pair}}\mathcal{L}_{\mathrm{pair}}+
    \lambda_{\mathrm{rank}}\mathcal{L}_{\mathrm{rank}}
    }.
    \label{eq:reliability_loss}
\end{equation}
We use $\lambda_{\mathrm{pair}}=0.2$,
$\lambda_{\mathrm{rank}}=0.1$, and $m_r=0.2$.

\subsection{Reliability-Guided Event Refinement}
\label{sec:event_refinement}

The event refiner processes the complete event voxel together with RGB
context and coarse depth:
\begin{equation}
    D_i^{\mathrm{raw}}=
    \mathcal{F}_{\mathrm{evt}}
    (\bm{V}_i,\bm{I}_i^{e_i},D_i^{c}).
\end{equation}
We deliberately do not replace $\bm{V}_i$ by
$\widehat{R}_i\odot\bm{V}_i$, because an imperfect reliability map could
irreversibly remove sparse event details before temporal encoding. Instead, we
express the event correction in log-depth space:
\begin{equation}
    \Delta_i^{\mathrm{raw}}(\bm{u})=
    \log\!\left(
    \frac{D_i^{\mathrm{raw}}(\bm{u})}
    {D_i^{c}(\bm{u})}
    \right).
    \label{eq:raw_log_residual}
\end{equation}

Reliability is calibrated using a threshold $\tau_r$ and temperature
$\gamma_r$:
\begin{equation}
    R_i^{\mathrm{cal}}(\bm{u})=
    \sigma\!\left(
    \frac{\widehat{R}_i(\bm{u})-\tau_r}{\gamma_r}
    \right).
\end{equation}
We combine it with a support floor $\rho$:
\begin{equation}
    S_i^{\mathrm{soft}}(\bm{u})=
    \rho+(1-\rho)S_i(\bm{u}),
\end{equation}
and optionally retain only the top-$\kappa$ fraction of calibrated
reliability values using the binary mask $\mathcal{M}_{\kappa}$. The resulting
reliability is
\begin{equation}
    \widetilde{R}_i(\bm{u})=
    R_i^{\mathrm{cal}}(\bm{u})
    S_i^{\mathrm{soft}}(\bm{u})
    \mathcal{M}_{\kappa}(\bm{u}).
\end{equation}
The soft output gate is
\begin{equation}
    A_i(\bm{u})=
    a_0+(1-a_0)\widetilde{R}_i(\bm{u}),
    \label{eq:soft_output_gate}
\end{equation}
where $a_0$ preserves a weak correction path for uncertain evidence.

The gated correction is bounded as
\begin{equation}
    \Delta_i(\bm{u})=
    \operatorname{clip}\!\left(
    \alpha_{\Delta}A_i(\bm{u})
    \Delta_i^{\mathrm{raw}}(\bm{u}),
    -\delta_{\max},\delta_{\max}
    \right),
    \label{eq:gated_log_residual}
\end{equation}
and the refined depth and point map are
\begin{align}
    D_i(\bm{u})
    &=D_i^{c}(\bm{u})\exp(\Delta_i(\bm{u})),
    \label{eq:refined_depth}\\
    \bm{P}_i(\bm{u})
    &=\bm{P}_i^{c}(\bm{u})
    \frac{D_i(\bm{u})}{D_i^{c}(\bm{u})}.
    \label{eq:refined_points}
\end{align}
Thus, all events contribute to feature extraction, but only reliable evidence
can produce a strong geometric correction.

To supervise the correction, we define a target in log-depth space:
\begin{equation}
    \Delta_i^{\ast}(\bm{u})=
    \operatorname{clip}\!\left(
    \log D_i^{\ast}(\bm{u})-
    \log D_i^{c}(\bm{u}),
    -\delta_t,\delta_t
    \right).
    \label{eq:residual_target}
\end{equation}
With reliability- and geometry-aware weight
\begin{equation}
    W_i^{\Delta}(\bm{u})=
    M_i(\bm{u})G_i(\bm{u})S_i(\bm{u})
    \left[r_f+(1-r_f)\widehat{R}_i(\bm{u})\right],
\end{equation}
the residual objective is
\begin{equation}
    \mathcal{L}_{\Delta}=
    \mathbb{E}_{W^{\Delta}}
    \left[|\Delta_i-\Delta_i^{\ast}|\right]
    +\lambda_{\nabla\Delta}
    \mathbb{E}_{W^{\Delta}}
    \left[\|\nabla\Delta_i-
    \nabla\Delta_i^{\ast}\|_1\right].
    \label{eq:residual_loss}
\end{equation}

We additionally discourage artificial corrections on geometrically planar
regions, including planar areas containing non-geometric events. Define
\begin{equation}
    W_i^{\mathrm{flat}}(\bm{u})=
    M_i(\bm{u})[1-G_i(\bm{u})].
\end{equation}
The corresponding protection loss is
\begin{equation}
    \mathcal{L}_{\mathrm{flat}}=
    \mathbb{E}_{W^{\mathrm{flat}}}
    \left[1-\langle
    \widehat{\bm{N}}_i,\bm{N}_i^{\ast}
angle\right]
    +\lambda_{0}\mathbb{E}_{M_i(1-S_i)}[|\Delta_i|].
    \label{eq:flat_protection}
\end{equation}
This term suppresses spurious normal bands without spatially filtering the
predicted residual and therefore preserves genuine high-frequency details.

\subsection{Overall Objective and Training Schedule}
\label{sec:training_objective}

The base geometry loss is
\begin{equation}
    \mathcal{L}_{\mathrm{base}}=
    \lambda_{d}\mathcal{L}_{d}+
    \lambda_{\mathrm{pts}}\mathcal{L}_{\mathrm{pts}}+
    \lambda_{\mathrm{pose}}\mathcal{L}_{\mathrm{pose}}.
\end{equation}
The geometry-refinement stage optimizes
\begin{equation}
    \boxed{
    \mathcal{L}_{\mathrm{geo}}=
    \mathcal{L}_{\mathrm{base}}+
    \mathcal{L}_{\mathrm{det}}+
    \lambda_{\Delta}\mathcal{L}_{\Delta}+
    \lambda_{\mathrm{flat}}\mathcal{L}_{\mathrm{flat}}+
    \lambda_{\mathrm{reg}}\mathcal{L}_{\mathrm{reg}}
    },
    \label{eq:geometry_objective}
\end{equation}
where $\mathcal{L}_{\mathrm{reg}}$ contains residual-magnitude, first-order,
second-order, and patch-periodicity regularization.

Training proceeds in three stages. First, the bounded token adapter is trained
with paired LDR observations and $\mathcal{L}_{\mathrm{tok}}$. Second, the
paired-token teacher is frozen, $R_i^{\ast}$ is exported, and ReliabilityNet is
trained with $\mathcal{L}_{R}$. Third, ReliabilityNet and the RGB backbone are
frozen while the depth/point heads and temporal event refiner are optimized
with $\mathcal{L}_{\mathrm{geo}}$. During inference, paired exposures,
ground-truth depth, and all target-construction cues are removed; the model
consumes only a single LDR sequence and its event voxels.
